\section*{第 \ref{ch:b1:biosyntheses-integration} 章——生物合成与整合的细胞}

\begin{solution}{exo:b1:biosyntheses-integration:1}
碳骨架（来自中心途径）、还原力（\omterm{def:b1:biosyntheses-integration:ppp}{NADPH}）、
能量（\omterm{def:b1:water-small-molecules:atp}{ATP}）。\omterm{def:b1:biosyntheses-integration:ppp}{戊糖磷酸途径}。
\end{solution}

\begin{solution}{exo:b1:biosyntheses-integration:2}
己糖激酶，由葡萄糖-6-磷酸酶绕过；磷酸果糖激酶，
由果糖-1,6-二磷酸酶绕过；丙酮酸激酶，由丙酮酸
羧化酶加 PEP 羧激酶绕过。
\end{solution}

\begin{solution}{exo:b1:biosyntheses-integration:3}
餐后约 1 小时达峰值 \qty{7.5}{mmol/L}，约 3 小时后回到
\qty{5}{}；早餐前降至最低，约 \qty{4.5}{mmol/L}，
即早餐前的最低点。
\end{solution}

\begin{solution}{exo:b1:biosyntheses-integration:4}
种子具有乙醛酸循环，它把两分子乙酰\omterm{def:b1:enzymes:enzyme}{辅酶} A 转化为一分子
四碳酸，不以 $\mathrm{CO_2}$ 形式损失碳，并由此
该酸生成糖。动物缺少这条捷径中的两种酶：
其乙酰\omterm{def:b1:enzymes:enzyme}{辅酶} A 进入\omterm{def:b1:respiration-fermentation:krebs}{克雷布斯循环}，两个碳以
$\mathrm{CO_2}$ 的形式离开，此时尚无净草酰乙酸生成。
\end{solution}

\begin{solution}{exo:b1:biosyntheses-integration:5}
每分子葡萄糖相当于 6 个 \omterm{def:b1:water-small-molecules:atp}{ATP}（两次羧化、两个 GTP、
磷酸甘油酸步骤中的两个 \omterm{def:b1:water-small-molecules:atp}{ATP}）。每个循环的净值：$2 - 6 = -4$ \omterm{def:b1:water-small-molecules:atp}{ATP}。
肝脏在冲刺之后借助氧气偿还——肌肉借用了
它无法及时以有氧方式合成的 \omterm{def:b1:water-small-molecules:atp}{ATP}。
\end{solution}

\begin{solution}{exo:b1:biosyntheses-integration:6}
每分子葡萄糖生成 2 分子乙酰\omterm{def:b1:enzymes:enzyme}{辅酶} A，需 4 分子葡萄糖。
经氧化支路，每分子葡萄糖生成 2 分子 \omterm{def:b1:biosyntheses-integration:ppp}{NADPH}，需 7 分子葡萄糖；合成一分子棕榈酸
需 11 分子葡萄糖，其中 4 分子提供碳，7 分子提供电子。
\end{solution}

\begin{solution}{exo:b1:biosyntheses-integration:7}
不同的酶使两个方向可以各自独立地开启或关闭；
不同的区室（合成在\omterm{def:b1:eukaryotic-cell:organelle}{细胞溶胶}中，
氧化在基质中）将两个中间产物库分隔开；\omterm{def:b1:biosyntheses-integration:ppp}{NADPH} 保持还原态以供
合成，\omterm{def:b1:respiration-fermentation:carriers}{NAD} 保持氧化态以供降解，因此
每个方向在各自的位置上都具有热力学优势。
丙二酰\omterm{def:b1:enzymes:enzyme}{辅酶} A 是合成途径中第一个定向中间产物，它抑制
将\omterm{def:b1:lipids:fattyacid}{脂肪酸}运入\omterm{def:b1:eukaryotic-cell:mitochondrion}{线粒体}的肉碱穿梭：
\omterm{def:b1:lipids:triglyceride}{脂肪}正在合成时，便无\omterm{def:b1:lipids:triglyceride}{脂肪}被燃烧。
\end{solution}

\begin{solution}{exo:b1:biosyntheses-integration:8}
丙氨酸 $+ \alpha$-酮戊二酸 $\rightleftharpoons$ 丙酮酸 $+$
\omterm{def:b1:biosyntheses-integration:nitrogen}{谷氨酸}。进食\omterm{def:b1:proteins:peptide}{蛋白质}后，过量的丙氨酸将其氮转给
\omterm{def:b1:biosyntheses-integration:nitrogen}{谷氨酸}（向右进行），氮以尿素形式排出，碳则用作燃料或
合成葡萄糖。禁食时，肌肉\omterm{def:b1:proteins:peptide}{蛋白质}被分解，其氨基
集中到丙酮酸上形成丙氨酸（在肌肉中向左进行），运往
肝脏，并在那里重新转化为丙酮酸和葡萄糖。
\end{solution}

\begin{solution}{exo:b1:biosyntheses-integration:9}
肝脏无法释放游离葡萄糖：两餐之间血糖会
危险地下降（低血糖）。葡萄糖-6-磷酸
积累，并被转入\omterm{def:b1:carbohydrates:polysaccharide}{糖原}，而肝脏无法将\omterm{def:b1:carbohydrates:polysaccharide}{糖原}
动员进入血液，肝脏因\omterm{def:b1:carbohydrates:polysaccharide}{糖原}而增大。过量的
葡萄糖-6-磷酸还会沿\omterm{def:b1:respiration-fermentation:glycolysis}{糖酵解}生成乳酸，使乳酸
在血液中升高。
\end{solution}

\begin{solution}{exo:b1:biosyntheses-integration:10}
$120/0.55 = \qty{218}{g}$ \omterm{def:b1:proteins:peptide}{蛋白质}每天，约 \qty{1.1}{kg} 肌肉。
若有酮体，则需 $40/0.55 = \qty{73}{g}$ \omterm{def:b1:proteins:peptide}{蛋白质}，
即 \qty{360}{g} 肌肉：损失减少三分之二，
依靠\omterm{def:b1:lipids:triglyceride}{脂肪}储备的存活时间由数周延长至数月。
\end{solution}

\begin{solution}{exo:b1:biosyntheses-integration:11}
同一磷酸化级联反应激活其中一个、使另一个
失活，因此开启一个阀门的信号同时关闭另一个。若两者
同时活跃，则每加入一个葡萄糖（消耗一个 UTP，相当于一个 \omterm{def:b1:water-small-molecules:atp}{ATP}）并被移除（游离），
每轮循环便耗费一个 \omterm{def:b1:water-small-molecules:atp}{ATP}，除热量外别无产物。
熊蜂恰恰在其飞行肌中运行着这样一个循环（以果糖-6-磷酸为\omterm{def:b1:enzymes:enzyme}{底物}），
在寒冷的清晨起飞前，将 \omterm{def:b1:water-small-molecules:atp}{ATP} 作为热源燃烧，
直至肌肉温度足以飞行。
\end{solution}

\begin{solution}{exo:b1:biosyntheses-integration:12}
每个交叉点上的酶只对周围的分子作出反应——
\omterm{def:b1:water-small-molecules:atp}{ATP}、AMP、柠檬酸、乙酰\omterm{def:b1:enzymes:enzyme}{辅酶} A、由激酶加上的磷酸基团——
并据此开启或关闭；两个方向由各自独立的酶控制，
这意味着开启一个方向可能会关闭另一个。
没有任何酶“知道”早餐或大脑；然而这些局部反应的总和，
使肝脏在进食后储存、夜间释放、
禁食时节约\omterm{def:b1:proteins:peptide}{蛋白质}：这是一种协调一致的策略，却无策划者。
所谓“计划”即阀门的连接方式，经选择而形成；
同时调控众多阀门的激素，使\omterm{def:b1:organism-environment:organism}{生物体}得以
作为一个整体行动。
\end{solution}

\begin{solution}{pb:b1:biosyntheses-integration:1}
\textbf{1.} 肝\omterm{def:b1:carbohydrates:polysaccharide}{糖原} $100\times 17 = \qty{1.7}{MJ}$，0.2 天；
肌\omterm{def:b1:carbohydrates:polysaccharide}{糖原} \qty{6.8}{MJ}，0.8 天；\omterm{def:b1:lipids:triglyceride}{脂肪} $12\,000\times 38 =
\qty{456}{MJ}$，54 天；\omterm{def:b1:proteins:peptide}{蛋白质} \qty{170}{MJ}，20 天（从不能
被完全利用）。
\textbf{2.} 肝\omterm{def:b1:carbohydrates:polysaccharide}{糖原}：肝脏具有葡萄糖-6-磷酸酶。
肌肉缺乏该酶，因此其\omterm{def:b1:carbohydrates:polysaccharide}{糖原}只能为自身
\omterm{def:b1:respiration-fermentation:glycolysis}{糖酵解}提供葡萄糖-6-磷酸。
\textbf{3.} $120 + 40 = \qty{160}{g}$ 每天。
\textbf{4.} $0.005\times 5\times 180 = \qty{4.5}{g}$：以 \qty{5}{g/h} 计，约 \qty{54}{min}。
\textbf{5.} $120\times 17 = \qty{2.0}{MJ}$，占 \qty{8.4}{MJ} 的四分之一。
\textbf{6.} 在肝脏中转化为\omterm{def:b1:lipids:fattyacid}{脂肪酸}和\omterm{def:b1:lipids:triglyceride}{甘油三酯}，
输出至\omterm{def:b1:lipids:triglyceride}{脂肪}\omterm{def:b1:body-plans-tissues:tissue}{组织}；胰岛素。
\textbf{7.} $160/24 = \qty{6.7}{g/h}$。
\textbf{8.} $100/6.7 = \qty{15}{h}$。
\textbf{9.} \omterm{def:b1:carbohydrates:polysaccharide}{糖原}提供三分之二，即 \qty{4.4}{g/h}：$100/4.4 =
\qty{22}{h}$。
\textbf{10.} 按第一种估计为次日 11:00，按第二种为 18:00——
两种情况下都需撑过整夜，余量
取决于\omterm{def:b1:biosyntheses-integration:gluconeogenesis}{糖异生}多早接管。
\textbf{11.} \qty{12}{h} 内的静息能量消耗：\qty{4.2}{MJ}；
\qty{60}{\%} 即 \qty{2.5}{MJ}；$2500/38 = \qty{66}{g}$ \omterm{def:b1:lipids:triglyceride}{脂肪}。
\textbf{12.} $160/0.55 = \qty{290}{g}$ \omterm{def:b1:proteins:peptide}{蛋白质}，即每天
\qty{1.45}{kg} 肌肉。
\textbf{13.} $160/180 = \qty{0.89}{mol}$ 葡萄糖；按每分子 6 个 \omterm{def:b1:water-small-molecules:atp}{ATP} 计
为 \qty{5.3}{mol} \omterm{def:b1:water-small-molecules:atp}{ATP}，以 \qty{50}{kJ/mol} 计
约 \qty{265}{kJ}——约占当日 \qty{8.4}{MJ} 的 \qty{3}{\%}：制造
葡萄糖并不昂贵，真正紧缺的是碳骨架，而非能量，
即身体所短缺的东西。
\textbf{14.} 非葡萄糖的能量消耗 $8.4 - 160\times 0.017 =
\qty{5.7}{MJ}$：$5700/38 = \qty{150}{g}$ \omterm{def:b1:lipids:triglyceride}{脂肪}，其甘油可提供
\qty{15}{g} 葡萄糖。
\textbf{15.} $(160 - 15)/0.55 = \qty{264}{g}$ \omterm{def:b1:proteins:peptide}{蛋白质}，即每天
\qty{1.3}{kg} 肌肉。
\textbf{16.} \qty{3}{kg} 肌肉\omterm{def:b1:proteins:peptide}{蛋白质}，按每天 \qty{264}{g} 计：
约 \qty{11}{d}。
\textbf{17.} $(80 - 15)/0.55 = \qty{118}{g}$ \omterm{def:b1:proteins:peptide}{蛋白质}，即每天
\qty{0.6}{kg} 肌肉；\qty{25}{d} 内损失一半肌肉——而在实际中
损失会进一步降低，降至每天 \qtyrange{20}{30}{g}，因为身体
会节约能量。
\textbf{18.} $456/8.4 = 54$ 天；以 \qty{6.5}{MJ} 计，为 70 天。
\textbf{19.} \omterm{def:b1:proteins:peptide}{蛋白质}并非储存物，而是机器本身：损失其三分之一
（呼吸肌、心脏、免疫蛋白）即会致死，而每天损失
\qty{100}{g} 或更多时，这会在数周内发生，而\omterm{def:b1:lipids:triglyceride}{脂肪}
可维持两个月。酮体使大脑能以\omterm{def:b1:lipids:triglyceride}{脂肪}为能源，减少
对葡萄糖的需求，从而将\omterm{def:b1:proteins:peptide}{蛋白质}损失降至三分之一，并将
存活时间延长至\omterm{def:b1:lipids:triglyceride}{脂肪}所设定的极限。
\textbf{20.} \omterm{def:b1:carbohydrates:polysaccharide}{糖原}磷酸化酶在 03:00 活跃，\omterm{def:b1:carbohydrates:polysaccharide}{糖原}合酶
在 09:00 活跃。
\textbf{21.} 胰高血糖素（03:00）与胰岛素（09:00）；两种酶均由
激酶级联反应磷酸化，再由磷酸酶使之去磷酸化。
\textbf{22.} 乙酰\omterm{def:b1:enzymes:enzyme}{辅酶} A 激活丙酮酸羧化酶；\omterm{def:b1:water-small-molecules:atp}{ATP}（及
柠檬酸）抑制磷酸果糖激酶。
\textbf{23.} 葡萄糖转化为丙酮酸生成 2 个 \omterm{def:b1:water-small-molecules:atp}{ATP}；丙酮酸转回
葡萄糖则需 6 个：每个往返循环白白耗去 4 个 \omterm{def:b1:water-small-molecules:atp}{ATP}，却无任何产物。三处旁路各有
独立的酶，同样的信号（\omterm{def:b1:water-small-molecules:atp}{ATP}、
乙酰\omterm{def:b1:enzymes:enzyme}{辅酶} A、磷酸化）激活其中一组、抑制另一组，
因此只有一个方向开放。
\textbf{24.} 没有胰岛素，阀门的设置与禁食时相同，无论血糖
水平如何：磷酸化酶、\omterm{def:b1:biosyntheses-integration:gluconeogenesis}{糖异生}和\omterm{def:b1:lipids:triglyceride}{脂肪}分解开启，
合成关闭；肝脏“认为”身体处于饥饿状态，于是
将葡萄糖和酮体倾入本已充满葡萄糖的血液，而依赖胰岛素的
\omterm{def:b1:body-plans-tissues:tissue}{组织}却无法摄取这些葡萄糖。
\textbf{25.} 无\omterm{def:b1:biosyntheses-integration:gluconeogenesis}{糖异生}时约 \qty{15}{h}，有\omterm{def:b1:biosyntheses-integration:gluconeogenesis}{糖异生}时约 \qty{22}{h}，
即一夜加一个上午；\omterm{def:b1:proteins:peptide}{蛋白质}损失最初约为每天 \qty{260}{g}，
待酮体供应大脑后降至 \qty{120}{g}（并进一步下降）；
\omterm{def:b1:lipids:triglyceride}{脂肪}储备可维持约 \qtyrange{55}{70}{d}。
\end{solution}
